\roteiro{2}

% Slide 6: cinco paradigmas (relatório §2.1, §2.2)
% As figuras 3 a 6 do relatório não entram: em cinco colunas ficariam ilegíveis.
\newcommand{\paradigma}[4]{%
  \begin{minipage}[t]{0.186\textwidth}
    \begin{painel}{#1}{#2}
      \raggedright #3\par\smallskip
      \destaque{Exemplo:} #4
    \end{painel}
  \end{minipage}}
\begin{frame}{Cinco paradigmas, cinco SGBDs}
  \begingroup\tcbset{equal height group=paradigmas}\small
  \paradigma{Relacional}{PostgreSQL}
    {Tabelas ligadas por chaves, com transações consistentes.}
    {pedidos e inventário.}\hfill
  \paradigma{Chave-valor}{Redis}
    {Cada chave aponta para um valor.}
    {carrinho e sessão.}\hfill
  \paradigma{Documento}{MongoDB}
    {Documentos JSON com campos variáveis.}
    {catálogo de produtos.}\hfill
  \paradigma{Colunar}{Cassandra}
    {Linhas com famílias de colunas.}
    {log de atividade.}\hfill
  \paradigma{Grafo}{Neo4j}
    {Nós e arestas com propriedades.}
    {recomendação de produtos.}\par
  \endgroup
  \smallskip
  Os quatro últimos são os bancos \destaque{NoSQL} (não relacionais).\par
  Um SGBD \destaque{multimodelo} reúne vários desses modelos no mesmo sistema.
\end{frame}

% Slide 7: trabalhos relacionados (relatório Cap. 3)
% Legenda: \tbok = o Cap. 3 afirma; \tbnd = o Cap. 3 não afirma.
\newcommand{\tbok}{\checkmark}
\newcommand{\tbnd}{\textemdash}
\begin{frame}{Trabalhos relacionados}
  \begingroup\footnotesize\bfseries\color{ufscazul} Comparação com os trabalhos relacionados\hfill
   \scriptsize\mdseries\color{ufsctexto}
   \tbok{} o Cap.~3 afirma \quad \tbnd{} o Cap.~3 não afirma\par\endgroup
  \vspace{1pt}
  \fontsize{8}{9}\selectfont % mesmo corpo do \scriptsize, entrelinha menor
  \setlength{\tabcolsep}{2.5pt}
  \setlength{\aboverulesep}{0.15ex}\setlength{\belowrulesep}{0.3ex}
  \renewcommand{\arraystretch}{1}
  \begin{tabular}{@{}E{0.235\textwidth}E{0.167\textwidth}E{0.205\textwidth}E{0.08\textwidth}%
      E{0.12\textwidth}E{0.115\textwidth}@{}}
    \toprule
    \bfseries Trabalho & \bfseries Foco & \bfseries Destinos & \bfseries Entrada CSV
    & \bfseries Mapeamento declarativo validado & \bfseries Revisão sem gravar \\
    \midrule
    Importadores nativos & importação & um SGBD por vez
      & \tbok\textsuperscript{1} & \tbnd & \tbnd \\
    \addlinespace[1.5pt]
    ArangoDB e \texttt{arangoimport} & importação & um SGBD multimodelo
      & \tbok & \tbnd & \tbnd \\
    \addlinespace[1.5pt]
    dbcrossbar & migração & PostgreSQL, BigQuery, nuvem & \tbok & \tbnd & \tbnd \\
    \addlinespace[1.5pt]
    Hackolade & modelagem visual & dezenas de alvos & \tbnd & \tbnd & \tbnd \\
    \addlinespace[1.5pt]
    Frameworks de acesso\textsuperscript{3} & acesso em tempo de execução & vários SGBDs
      & \tbnd & \tbnd & \tbnd \\
    \addlinespace[1.5pt]
    \textit{Polystores}\textsuperscript{2} (BigDAWG) & consulta federada & fontes heterogêneas
      & \tbnd & \tbnd & \tbnd \\
    \midrule
    \onslide<2->{\destaque{Polyglot Import CSV}} & \onslide<2->{\destaque{importação}}
      & \onslide<2->{\destaque{cinco SGBDs, cinco paradigmas}}
      & \onslide<2->{\destaque{\tbok}}
      & \onslide<2->{\destaque{\tbok{} (JSON Schema)}}
      & \onslide<2->{\destaque{\tbok{} (\texttt{-{}-dry-run})}} \\
    \bottomrule
  \end{tabular}
  \par\vspace{2pt}
  {\fontsize{8}{9}\selectfont\color{black!70}
   \textsuperscript{1} Redis: arquivo em protocolo bruto. \quad
   \textsuperscript{2} Uma só consulta sobre vários SGBDs (consulta federada).\par
   \textsuperscript{3} Interface para ler e gravar em vários SGBDs (PEREIRA; GUERRA; ROSA, 2025;
   SACHDEVA; VASAVA, 2024).\par}
\end{frame}
